% =====================================================================
\section{Discussion}\label{sec:discussion}
% =====================================================================
This section interprets the held-out outcome against the model's conditional guarantee, draws its consequences for a site and states its limits.

\subsection{What the held-out comparison shows}\label{subsec:interpretation}
The article asked whether a reserved margin, historical scenarios with activation, or both keep a visit-minimising layout inside a two-sided share band on later orders. At the held-out deployment the outcomes divided along the margin factor. Both arms that reserved part of the band satisfied it on every optimiser seed, TIGHT meeting the predeclared primary endpoint, and neither arm without a reserve satisfied it on any seed. Historical scenarios with activation did not change which arms satisfied the band, although they changed how far station shares moved inside it. The factorial was designed for this separation. In the exploratory campaigns TIGHT differed from HIST+ACT both in its margin and in its lack of scenarios, so the two factors could not be told apart there; the held-out design crossed them at a later origin. What the design does not do is hold layouts fixed. The arms are different time-capped solves, and the pattern describes one deployment under equal configured solve budgets; it does not isolate the margin from solver and layout effects.

The training diagnostics suggest a reading of the pattern without establishing it. Every layout returned within the time limit used essentially its whole training allowance, the behaviour anticipated in Section~\ref{sec:introduction} for an optimiser rewarded only for fewer visits. A layout optimised at the full band then enters the future window with at least one station on an edge of the band on its training scenarios, whereas a layout optimised inside the tightened band enters it with the reserved margin still available on those scenarios. On this window that difference coincided with the difference between satisfying and missing the band. It is a hypothesis about the outcome, not its demonstrated cause; Section~\ref{subsec:limitations} states it with its conditions and the experiment that would test it.

\subsection{Conditional robust feasibility and observed compliance}\label{subsec:feasibility}
Conditional robust feasibility and observed future compliance are different quantities, and the held-out window separates them. If a layout satisfies the robust rows and the future product mix lies in the declared set, every station respects its band; the model says nothing about whether the future lies in the set and assigns no probability to compliance. On this window the premise failed for every returned layout, each of which had a station whose realised share fell below the lowest share its modelled set allowed.

This departure is a demonstrated boundary condition of the conditional guarantee: the assumption is named, and the evidence shows where it fails. The historical hull at the held-out origin is spanned by the mixes of eleven blocks and of the pooled history, a set of low affine dimension relative to the number of products, and activation widens it only towards products that had no historical workload; the study records that a later window fell outside such a set for every returned layout at this origin. The finding concerns historical block scenarios with this activation budget at this site and is not evidence about robust optimisation in general; other uncertainty families were not tested. Nor does it bear on the probabilistic guarantees of budgeted, sample-based or distributionally robust methods \citep{bertsimas2004price,luedtke2008sample,vanparys2021data}, which rest on random-sampling or independence assumptions that consecutive blocks of one order stream are not shown to meet (Section~\ref{subsec:rw_robust}).

The sufficient condition of the second proposition did not hold either. Its premise, a product-mix distance from pooled history within the reserved margin, failed on every analysed window, the held-out distance being many times the reserve (Section~\ref{subsec:descriptive}). That condition is a certificate that can hold on a given window, not a rule for choosing a reserve, and its failure does not point to any particular alternative rule. The evidence is consistent with product-level changes cancelling within stations, since each station's share aggregates the workload of many products, but the study does not measure that cancellation. The held-out result is a positive margin finding whose mechanism is not established.

Within its scope the study excludes three readings of the held-out outcome. (i) The compliance of the margin arms did not require the future product mix to lie in their modelled sets, because the mix lay outside them. (ii) The band held although the reserve was smaller than the incumbent's largest historical block deviation (exploratory history, incumbent layout), so a reserve at least that large was not needed for compliance at this deployment. (iii) The product-level total-variation condition does not account for compliance on these windows, because it did not hold.

\subsection{Staying close to historical shares versus visits}\label{subsec:tradeoff}
Section~\ref{subsec:heldout} shows that historical scenarios with activation kept the largest station deviation lower at higher mean visits per order, without changing which arms satisfied the band. The two margin arms therefore offer a site two readings of the same band. If only compliance with the band matters, TIGHT satisfied it on every seed at the lower mean visit count. If the site also values how far station shares move inside the band, HIST+ACT-T held the largest deviation lower at visits per order that were higher on every seed. Which reading applies is a site decision that the study does not make.

The trade-off is descriptive. Scenarios and activation are bundled, so the change in profile cannot be assigned to either, and the largest deviation is attained at one station per layout, so a lower value does not show that a whole station profile is closer to its targets.

Every optimised arm, the two margin arms included, made fewer future visits per order than the incumbent on the same window. Because the targets are the incumbent's own historical shares (Section~\ref{subsec:heldout}), the incumbent's compliance and the misses of the arms without a margin are outcomes of layouts that started from different positions relative to the band, not a ranking of the layouts' protection.

\subsection{Operational meaning for a site}\label{subsec:operational}
For a site that wants fewer station visits while keeping every station's workload share inside a band, the policy tested here is to optimise inside a band tightened so that part of the policy band stays in reserve. Reserving part of a constraint during optimisation is related to the protection levels of budgeted robust optimisation (Section~\ref{subsec:rw_robust}); the study adds a held-out test of a fixed reserve for share targets in one pick-and-pass warehouse.

Checking the reserve against historical station-level variation did not predict compliance here (Section~\ref{subsec:descriptive}): the reserve was smaller than the incumbent layout's largest single-block share deviation on the earlier exploratory history, yet both margin arms satisfied the band. The check therefore did not validate the reserve, and because that variation belongs to the incumbent's own layout, it informs but does not set the reserve for an optimised layout. The reserve used here is a policy tested on one case that still needs validation, not a calibrated or automatic margin rule.

The design suggests three practices, none of which was tested as an operating procedure. First, targets and the reserve can be fixed and recorded at the deployment origin, before later orders are scored, as this study did, so that neither is fitted to the orders on which it is judged. Second, realised station shares are worth following on both sides of the band (Section~\ref{subsec:uncertainty} explains why caps alone do not protect floors): HIST+ACT breached the floor on every held-out seed, NOM had four or five stations below it on every seed (Table~\ref{tab:heldout_seeds}), and the exploratory misses of the scenario arms were mostly on the floor side (Section~\ref{subsec:exploratory}). Third, the minimum required slack of a candidate layout shows how much of the band it uses on its training history: about $\delta$ without a reserve, leaving none, and about $(1-\lambda)\delta$ with one, leaving $\lambda\delta$, so it confirms that the reserve is intact on the training history. Every returned layout used its whole training allowance, whether or not it complied, so the value does not indicate whether a layout will comply; the values here describe time-capped returned layouts. A site with order dates may prefer calendar horizons, which this undated protocol defers rather than rules out.

\subsection{Limitations}\label{subsec:limitations}
The held-out evidence is one industrial warehouse, one deployment origin, one horizon of $|P|$ complete orders and three optimiser seeds under one frozen policy (Section~\ref{subsec:scope}). The seeds describe optimiser variability, not variation in future demand, and nothing is claimed about other origins, horizons or warehouses. Each limitation below states what failed or remains unknown, under which conditions, whether its explanation is demonstrated or hypothetical, and which experiment would resolve it; no experiment is run for this article. Negative and mixed results are reported as found, and no band, activation budget, arm, instance or horizon was changed or removed to improve a table. The predeclared primary endpoint, whose failure would have closed the fixed-policy claim, was met; the two arms without a reserve missed the band on every seed (Section~\ref{subsec:heldout}), and the misses of HIST+ACT are discussed below.

\paragraph{Demonstrated boundary conditions.} The conditional feasibility of the model assumes that the future product mix lies in the modelled set. That assumption failed for every returned held-out layout, at the held-out origin with the historical block hull and the activation budget used there, and the failure is demonstrated, because a station share outside its modelled range proves that the future lay outside the set. Whether another uncertainty family would contain later windows while leaving layouts with useful visit counts is unknown; a predeclared comparison of uncertainty families at further origins would answer it. The sufficient product-mix condition for the reserve assumes a distance from pooled history within the reserved margin, and every analysed window exceeded it, which is also demonstrated. Because the condition is sufficient only, its failure leaves open which property of the windows allowed the margin arms to comply; measuring the within-station cancellation of product-level change on scored windows would address that question.

\paragraph{Empirical limitations.} The held-out misses of HIST+ACT are an empirical limitation of scenario-and-activation protection without a reserved margin. Under the frozen policy the arm satisfied the two-sided band on none of the three seeds, breaching below the floor on every seed and above the cap on two (Table~\ref{tab:heldout_seeds}). The conditions were one origin, one horizon of $n=|P|$ orders, $\delta=0.02$ and $\nu=0.01$, with scenarios and activation bundled and every solve time-capped. One possible explanation is that each returned layout used essentially its whole training allowance while the future fell below its modelled set at several stations (Table~\ref{tab:heldout_diag}), so that a station at the edge of the band on its training scenarios could move outside it. That explanation is hypothetical: HIST+ACT-T also used its whole, tightened allowance and also had stations below its set, yet satisfied the band on every seed. Further held-out origins with a predeclared $\lambda$ frontier for both scenario arms, together with layouts solved closer to proven optimality, would resolve it; neither is run for this article. Separately, and only at the exploratory origin, a wider band did not bring the scenario arms into compliance: with $\delta=0.03$ neither HIST nor HIST+ACT satisfied the band at any horizon (Table~\ref{tab:exploratory}). Each value of $\delta$ was a different model solved with one seed, the mechanism is not isolated, and a predeclared band-width and margin response at held-out origins would be needed to explain it. This item is exploratory and is not pooled with the held-out one.

\paragraph{Experimental limitations of the design.} The design cannot separate historical scenarios from activation, which enter the factorial as one factor; a factorial that also crosses activation, for instance by running HIST with and without the reserved margin beside the four arms, would separate them. It cannot isolate the reserved margin from solver and layout effects, because every arm is a different time-capped solve; repeated deployments with layouts solved closer to proven optimality would reduce that confounding. It has one held-out origin, one horizon and one warehouse, and it ran no $\lambda$, $\delta$ or $\nu$ frontier at that origin, so no validated rule for choosing the reserve was obtained. The history of the held-out origin also contains the exploratory origin's data, so the held-out test concerns the same order stream that shaped the policy. Predeclared deployments on later orders or at other sites, each with a margin frontier fixed in advance, would address these limits. The study examined a single band shape, absolute in share units, with targets defined as the incumbent's own shares; a band relative to each target and targets defined in other ways were not studied, and findings about this band do not transfer to them.

\paragraph{Experimental limitations of the data.} The industrial snapshot (catalogue, incumbent layout, slot counts, freeze mask and order-retention rule) was reconstructed from the whole export and assumed known at every origin, so the experiment is retrospective, and prefix-only estimation does not show that this metadata was available at each origin; a prospective deployment with metadata frozen at the origin would remove that qualification. Chronology rests on order identifiers by assumption, and the dated alternative is deferred, not disproved; the same comparison on a dated extract would test it. The catalogue is closed, so products that appear later are outside scope, and workload counts product-order pairs rather than processing time. The like-for-like drift and dispersion values take each block's mass fraction by order count, because block line mass is not stored. The synthetic instances come from a stationary generator and inform only finite-sample and solver behaviour, not robustness to drift.

\paragraph{Implementation limitations.} Every held-out layout is the best the solver returned within its time limit, with a relative gap of about 98.5\% on the training visit objective, so no visit figure is an optimal value and the visit difference between arms is not identified exactly. Which layout a time-capped solve returns, and therefore its future compliance, may change with solve effort as well as with the seed, beyond what three seeds at one origin show. Runs that returned no layout within their time limit are not evidence of infeasibility, and no infeasibility certificate was obtained. Six earlier diagnostic solves of the exploratory stage, which bounded the minimum required slack, were logged without saving their layouts, so their bounds are logged results only; they are distinct from the minimum required slack of the scored layouts reported in Section~\ref{subsec:heldout}, and no result in this article uses them. Longer solves, or solves with informative bounds, of the held-out models would resolve the first of these limitations.

\subsection{What the study does not establish}\label{subsec:rulesout}
The study does not establish that a reserved margin keeps the band at other origins, horizons, warehouses or reserve sizes, or that the margin produced the compliance observed. It does not identify the separate contributions of historical scenarios and of activation, the visit difference that exact optimisation with and without protection would give, a rule for choosing $\lambda$ or $\delta$, or the mechanism by which the margin arms complied outside their modelled sets. Nor does it show anything about robustness to drift from the synthetic instances, that dates are unnecessary for operational decisions in general, or that robust optimisation over other uncertainty families would behave as the historical block hull did here.
